% Regression: t2_pentagon4 — renders ^2F ^3F = ^3F ^2F ^2F — a horizontal MPO string (red.
% Target: RMP 2011.12127 fig3_pentagon4.pdf (pentagon-equation component).
% Formula:  ^2F ^3F = ^3F ^2F ^2F  —  a horizontal MPO string (red,
% directed east) slides through a physical splitting vertex (gray disc,
% root arrow in from below, two legs out the top):
%   LHS: the string crosses the root leg BELOW the vertex;
%   RHS: the string crosses BOTH split legs ABOVE the vertex.
% Genuinely non-grid (free crossings + a trivalent vertex): tenkzfree.
% Honest gaps: \tnjoin has no wire-hue key, so the red (MPO) vs black
% (physical) distinction is carried only by geometry; the splitting
% vertex renders as a house bead, not a gray disc.
\documentclass[varwidth=19cm, border=6pt]{standalone}
\usepackage{amsmath, amssymb}
\usepackage{tenkz}
\begin{document}

\[
  \begin{tenkzfree}
    % the splitting vertex: root in from below, two legs out the top
    \tnput[dot, ports={south:physical}]{V}{(0,0)}{}
    \tnjoin[dir]{0mm,-16mm}{V.south}
    \tnjoin[dir, route=curve, out=130, in=-90]{V.north west}{-8mm,14mm}
    \tnjoin[dir, route=curve, out=50, in=-90]{V.north east}{8mm,14mm}
    % the MPO string, below the vertex, crossing the root leg
    \tnjoin[dir]{-9mm,-9mm}{9mm,-9mm}
  \end{tenkzfree}
  \;=\;
  \begin{tenkzfree}
    \tnput[dot, ports={south:physical}]{W}{(0,0)}{}
    \tnjoin[dir]{0mm,-12mm}{W.south}
    \tnjoin[dir, route=curve, out=130, in=-90]{W.north west}{-8mm,18mm}
    \tnjoin[dir, route=curve, out=50, in=-90]{W.north east}{8mm,18mm}
    % the string, above the vertex, crossing BOTH split legs
    \tnjoin[dir]{-13mm,11mm}{13mm,11mm}
  \end{tenkzfree}
\]
\[
  {}^{2}F\,{}^{3}F \;=\; {}^{3}F\,{}^{2}F\,{}^{2}F
\]

\end{document}
